\documentclass[11pt]{article}

\usepackage[margin=25mm]{geometry}

\usepackage{amsmath,amssymb,fontspec,unicode-math}

\IfFontExistsTF{DejaVu Serif}{\setmainfont{DejaVu Serif}}{\setmainfont{Latin Modern Roman}}

\IfFontExistsTF{Latin Modern Math}{\setmathfont{Latin Modern Math}}{\setmathfont{DejaVu Math TeX Gyre}}

\usepackage{xurl,hyperref}

\setlength{\parindent}{0pt}\setlength{\parskip}{6pt}

\begin{document}

\title{Act Only When It Pays: Efficient Reinforcement Learning for LLM Reasoning via Selective Rollouts}

\author{ELVIA preliminary AI analysis}\date{21 October 2026}\maketitle

Preliminary assessment: Minor.\par

Independent scientific review remains pending. Author adoption has not been requested.

The release can finish adaptive sampling below its effective-batch target and uses mean-threshold histories rather than the written zero-variance streak. Both findings are bounded implementation discrepancies; historical training and efficiency effects are unresolved.

\section*{Review scope}

Full supplied technical report read alongside the corresponding original-paper claim passages and applicable author-release source. This pass independently checks the stated reasoning and bounds; historical executions were not rerun.

\section*{Problem 1: Adaptive refill can exit with too few effective prompts}

Source: Algorithm 1 and Appendix F; verl/trainer/ppo/ray\_trainer.py, lines 1141–1182

Assessment: supported; impact: minor

\subsection*{Source claim (paraphrase)}

The algorithm promises a full effective batch, but last\_sampling bypasses the actual accumulated-count condition after a predicted final refill.

\subsection*{Analysis}

Use the documented \(256\)-prompt target, raw cap \(384\), and \(8\) responses per prompt. A first raw batch yields \(192\) effective groups; the estimated refill is \(1.25(256-192)/(1/2)=160\), setting last\_sampling. If that refill yields \(48\), the source exits with \(240<256\) groups. Mixed \(0.1/1\) rewards and constant \(0.1\) rewards realize the effective/ineffective groups. No premise forces identical realized yields across batches. With no effective groups, the code first divides by \(1-\alpha=0\), before its cap can apply.

\subsection*{Finding and scope}

The active release can violate the effective-batch invariant or fail on an admissible zero-yield batch. Historical occurrence and resulting speed/accuracy changes remain unknown.

\subsection*{Proposed correction}

Permit successful exit only when the accumulated count meets the target; explicitly branch on zero yield and use a retry/abort policy. This proves successful-exit correctness, not termination when effective prompts never occur.

The described correction remains a proposal; author adoption or a corresponding erratum is not established.

\section*{Problem 2: Mean-threshold runs differ from zero-variance streaks}

Source: Equations (3)–(5); presampling\_selector.py, easy\_probabilistic\_skip and hard\_probabilistic\_skip; ray\_trainer.py, lines 1121–1132

Assessment: supported; impact: minor

\subsection*{Source claim (paraphrase)}

The written retention rule uses the most recent consecutive zero-variance groups. The release stores means, counts separate easy/hard threshold runs and applies a retention floor.

\subsection*{Analysis}

An all-\(1\) group followed by an all-\(0.1\) group has equality streak \(z=2\), hence retention \(0.5^2=0.25\); the hard mean-threshold run counts only the latest group and retains with \(0.5\). A group with four \(0\) and four \(0.1\) rewards has mean \(0.05\), variance \(0.0025>0\), so the written rule retains with \(1\), but the hard branch can retain with only \(0.5\). The rewards are allowed by the released formatting/correctness scheme. Constant same-type streaks agree, and a strictly binary reward domain would avoid the second example; the actual multivalued domain does not.

\subsection*{Finding and scope}

The released pre-filter is a different history statistic, not an implementation of the equality indicator. No claim is made that every skipped prompt would improve future training.

\subsection*{Proposed correction}

Record the equality indicator before reducing groups to means and update the trailing streak by increment/reset; alternatively describe the released threshold runs. Disclose the \(0.01\) retention floor separately from the minimum exploration base.

The described correction remains a proposal; author adoption or a corresponding erratum is not established.

\section*{Limits}

Preliminary AI source/argument assessment; independent human scientific adjudication remains pending.\par

No author-provided or retained verification program was executed. Published empirical results were not independently reproduced.\par

\section*{Source identity}

NeurIPS 2025 official proceedings PDF (exact retained local copy inspected)

\url{https://proceedings.neurips.cc/paper_files/paper/2025/file/b42fb82fc3ffb88b872c4714f093875d-Paper-Conference.pdf}

Reviewed PDF SHA-256: \nolinkurl{8db4dc7d94f2f03231aeb53546faf32bffd494c7ca06625fb4f7ab85f1bacc45}

The SHA-256 identifies the exact retained PDF inspected in this rapid source/argument check. Original supplied report bytes are retained locally. The original PDF retrieval timestamp is taken from matching retained metadata.

\end{document}
